\chapter{LightGBM for Tabular Volatility}
\label{ch:lightgbm}

LightGBM \citep{Ke2017} is the primary ML model in the pipeline.
It ingests the full engineered feature set from Layers~0--7 and optimizes a custom $\QLIKE$ objective.

%% ──────────────────────────────────────────────
\section{What Goes In}
\label{sec:lgbm-features}
%% ──────────────────────────────────────────────

Table~\ref{tab:feature-layers} summarizes the raw feature counts per layer before expansion.
After applying \{level, change, z-score\} transformations to each base feature, the final input dimensionality is approximately 80--120.

\begin{table}[htbp]
\centering
\caption{Feature layers feeding the LightGBM model.}
\label{tab:feature-layers}
\small
\begin{tabular}{@{}clr@{}}
\toprule
\textbf{Layer} & \textbf{Features} & \textbf{Count} \\
\midrule
0 & $\log \RV$ (d/w/m), RQ, RQ interaction  & 5 \\
1 & $\RV^+$, $\RV^-$, signed jumps, $C$, $J$ & 6 \\
2 & ATM IV, $\VRP$, skew, term structure, $\VVIX$, butterfly, IV--RV gap, event IV & 9 \\
3 & Price acceleration, OBI, spread, VPIN, Kyle $\lambda$ & 9 \\
4 & Treasury slope, FX vol, credit spread, DY index & 8 \\
5 & Calendar: FOMC, OpEx, month/weekday dummies & $\sim$15 \\
6 & Memory: Hurst exponent, fractional $d$, ACF features & $\sim$10 \\
7 & Sentiment: FinBERT scores, attention proxies & $\sim$15--55 \\
\midrule
  & \textbf{Raw total} & $\sim$37--57 \\
  & \textbf{After \{level, change, z-score\} expansion} & $\sim$\textbf{80--120} \\
\bottomrule
\end{tabular}
\end{table}


%% ──────────────────────────────────────────────
\section{Configuration}
\label{sec:lgbm-config}
%% ──────────────────────────────────────────────

Table~\ref{tab:lgbm-config} gives the reference configuration, adapted from the Optiver Kaggle competition (91st-percentile solution).

\begin{table}[htbp]
\centering
\caption{LightGBM reference configuration.}
\label{tab:lgbm-config}
\small
\begin{tabular}{@{}llp{5.5cm}@{}}
\toprule
\textbf{Parameter} & \textbf{Value} & \textbf{Notes} \\
\midrule
\texttt{learning\_rate}        & 0.05    & \\
\texttt{num\_leaves}           & 255     & \\
\texttt{min\_data\_in\_leaf}   & 255     & Prevents overfitting on rare regimes \\
\texttt{n\_estimators}         & 10{,}000 & With early stopping \\
\texttt{early\_stopping\_rounds} & 400   & \\
\texttt{boosting\_type}        & \texttt{dart} & Dropout regularization \\
\texttt{objective}             & Custom  & $\QLIKE$ loss; not a built-in objective \\
\texttt{metric}                & Custom  & $\QLIKE$ evaluation for early stopping \\
\bottomrule
\end{tabular}
\end{table}


%% ──────────────────────────────────────────────
\section{Custom QLIKE Objective}
\label{sec:qlike-objective}
%% ──────────────────────────────────────────────

$\QLIKE$ is not a built-in LightGBM loss.
Training requires a custom objective function that returns the gradient and Hessian of the $\QLIKE$ loss with respect to the predicted value, plus a separate custom evaluation metric for early stopping.
$\QLIKE$ is robust to noise in the volatility proxy: it ranks forecasts consistently even when $\RV$ is measured with error, unlike MSE \citep{Patton2011}.
All training and prediction operate in $\log$-$\RV$ space; convert back to levels only for final $\QLIKE$ evaluation.


%% ──────────────────────────────────────────────
\section{SHAP Interpretability}
\label{sec:shap}
%% ──────────────────────────────────────────────

$\SHAP$ values \citep{Lundberg2017} provide feature importance and interaction effects for each prediction.
They are required for the GS presentation and for defending model behavior under scrutiny.
However, single-model importance measures (gain, permutation, $\SHAP$) are unstable across refits when features are near-substitutes (e.g., VIX, $\VVIX$, ATM IV all proxy for the same latent factor).
This instability motivates the Rashomon analysis in Chapter~\ref{ch:rashomon}, which quantifies importance across the set of near-optimal models rather than relying on a single fit.


%% ──────────────────────────────────────────────
%% Boxes
%% ──────────────────────────────────────────────

\begin{warning}[Baseline First]
The fitting scheme for the $\HAR$ baseline matters more than the choice of ML model \citep{HARdToBeat2024}.
A properly fitted $\HAR$ using OLS with Newey--West standard errors and an expanding or rolling window is essential before claiming ML improvement.
Without this baseline, apparent ML gains may reflect a weak benchmark, not genuine forecasting ability.
\end{warning}

\begin{keyidea}[Why LightGBM]
Gradient-boosted trees dominate tabular volatility forecasting benchmarks \citep{ChristensenSiggaardVeliyev2023, BrancoRubesamZevallos2024}.
They handle mixed feature types (continuous RV lags, categorical calendar dummies, ordinal sentiment scores) without preprocessing, capture nonlinear interactions automatically, and train in seconds on datasets of this size ($\sim$5{,}000 rows $\times$ 120 features).
\end{keyidea}
